\documentclass[11pt]{article}
\usepackage[margin=1.1in]{geometry}
\usepackage{amsmath,amssymb,amsthm}
\usepackage{array}

\newtheorem{theorem}{Theorem}
\newtheorem{proposition}[theorem]{Proposition}
\newtheorem{lemma}[theorem]{Lemma}
\newtheorem{corollary}[theorem]{Corollary}
\newtheorem{conjecture}[theorem]{Conjecture}
\newtheorem{openlemma}[theorem]{Open Lemma}
\theoremstyle{definition}
\newtheorem{definition}[theorem]{Definition}
\theoremstyle{remark}
\newtheorem{remark}[theorem]{Remark}

\newcommand{\Z}{\mathbb{Z}}
\newcommand{\R}{\mathbb{R}}
\newcommand{\E}{\mathbb{E}}
\newcommand{\CV}{\mathrm{CV}}
\newcommand{\Var}{\mathrm{Var}}
\newcommand{\cov}{\mathrm{cov}}
\newcommand{\lam}{\lambda}
\newcommand{\gam}{\gamma}
\newcommand{\bet}{\beta}
\newcommand{\kap}{\kappa}
\newcommand{\dl}{\delta}

\title{The attenuation constant of the $3x+1$ cascade:\\
an exact decomposition and a stability program}
\author{Martien de Jong\\[3pt]
{\normalsize with computational collaboration: Jengo (AI)}}
\date{July 2026}

\begin{document}
\maketitle

\begin{abstract}
The full-density program for the $3x+1$ problem reduces, by the companion
papers, to one Open Lemma: the attenuation constant $\kap$ of the
Krasikov--Lagarias minimum is uniformly bounded away from $1$
($\kap_\infty = 0.839 \pm 0.002$ measured, nothing proved). This paper is
the proof-attempt manuscript for that lemma. The engine is a one-line exact
formula: writing the member increments of a sibling triple as
$\Delta_i = t + \dl_i$ (common move plus centered fine moves) and the
static within-triple gaps as $g_i$, the increment of the triple minimum is
\emph{exactly} $t + S$ with $S = \min_i(\dl_i + g_i)$. The common move
therefore passes with slope exactly $1$ --- recovering the directional
low-pass --- and \emph{all} attenuation is carried by the scalar $S$.
Squaring gives an exact decomposition,
$\kap^2 = [\Var(t) + 2\cov(t,S) + \Var(S)]\,/\,[\Var(t) + \Var(\dl)]$,
verified to machine precision on certificate data
($0.7901 = 0.7901$ at $P = 4$; $0.7431$ at $P = 6$, $k = 13$), so that
$\kap < 1$ if and only if $2\cov(t,S) + \Var(S) < \Var(\dl)$. Both
sufficient mechanisms are measured with the right sign and size:
$\cov(t,S) < 0$ (selection-weighted mean reversion) and
$\Var(S)/\Var(\dl) = 0.639$, $0.499 < 1$ (min-variance contraction under
intense competition: the selected member is non-argmin with probability
$0.658$--$0.665 \approx 2/3$, and the median second gap is only
$0.13\,\sigma_\dl$). We prove a clean quantitative min-variance contraction
lemma --- an exact identity with no hypotheses, plus a contraction bound
whose symmetry hypothesis we state explicitly and confront with the data
--- and assemble the conditional theorem: \emph{if} the mean-reversion
sign and the competition condition hold uniformly in depth, \emph{then}
$\kap \leq \kap_{\max} < 1$ uniformly, the Open Lemma closes, and by the
companion chain $\gam_k \to 1$: density $x^{1-\varepsilon}$ for every
$\varepsilon > 0$. What remains unproved is exactly the two hypotheses;
we state the proof obligations, the supporting measurements at every
available depth, and the ways the program could fail.
\end{abstract}

\section{Introduction}
\label{sec:intro}

Let $T(n) = n/2$ for even $n$, $T(n) = (3n+1)/2$ for odd $n$, and let
$\pi_1(x)$ count the $n \leq x$ whose $T$-orbit reaches $1$. The
difference-inequality method of Krasikov and Lagarias \cite{KL03} converts
a finite feasibility certificate at $3$-adic depth $k$ into a density
bound $\pi_1(x) \geq x^{\gam_k}$. The companion papers \cite{Temp26,
Drift26, Gamma26} assemble a program that would push $\gam_k \to 1$: the
tempering law and the edge-rate theorem $d\gam/dq = 1/\ln(4/3)$
\cite{Temp26}, the mass-conservation, balance, and directional low-pass
identities \cite{Drift26}, and the synthesis as a single conditional
theorem \cite{Gamma26}. Every link of that chain is proved except one, the
\emph{Open Lemma}: the attenuation constant $\kap(P;k)$ of the
Krasikov--Lagarias minimum --- the ratio of the standard deviation of the
triple-minimum increment to that of the member increments, under one
application of the edge-linearized operator --- satisfies
$\kap(P;k) \leq \kap_{\max} < 1$ uniformly in scale $P$ and depth $k$.
Measured: $\kap < 0.95$ everywhere and $\kap_\infty = 0.839 \pm 0.002$,
flat across $k = 13, 15, 17, 19$ (Prop.~35 of \cite{NOTE}); proved:
nothing.

This paper is the proof attempt. Its contributions are:

\begin{enumerate}
\item \textbf{The $S$-formula} (Lemma~\ref{lem:S}; proved, one line). In
the natural coordinates --- common move $t$, centered fine moves $\dl_i$,
static gaps $g_i$ --- the increment of the triple minimum is exactly
$t + S$ with $S = \min_i(\dl_i + g_i)$. The common move passes with slope
exactly $1$ (recovering the directional low-pass, Prop.~34b of
\cite{NOTE}, as a corollary), and the entire attenuation mechanism is
compressed into the single scalar $S$.
\item \textbf{The exact $\kap$ decomposition}
(Corollary~\ref{cor:decomp}; proved, and verified to machine precision on
certificate data):
\[
\kap^2 = \big[\Var(t) + 2\cov(t,S) + \Var(S)\big] \,/\,
\big[\Var(t) + \Var(\dl)\big].
\]
Hence the Open Lemma is \emph{equivalent} to a two-term inequality:
$2\cov(t,S) + \Var(S) < \Var(\dl)$, uniformly.
\item \textbf{A min-variance contraction lemma}
(Lemma~\ref{lem:contraction}; proved). An exact identity --- valid with
no hypotheses whatsoever --- expresses the variance deficit of a minimum
as a product of one-sided swap means, $2\,\E[W^+]\,\E[W^-]$, where $W$ is the difference of the competing
candidates: the min contracts variance exactly when the \emph{order
swaps} with positive probability in both directions. A quantitative
contraction bound follows under an explicitly stated reflection-symmetry
hypothesis; we say precisely which hypotheses the certificate data
supports and which are untested.
\item \textbf{The conditional stability theorem}
(Theorem~\ref{thm:main}; proved as stated). If the mean-reversion sign
$\cov(t,S) \leq 0$ and the competition condition hold uniformly in depth,
then $\kap \leq \kap_{\max} < 1$ uniformly: the Open Lemma closes, and by
\cite{Gamma26} $\gam_k \to 1$, i.e.\ $\pi_1(x) \geq x^{1-\varepsilon}$
for every $\varepsilon > 0$.
\item \textbf{The honest gap} (Section~\ref{sec:gap}). What remains is to
prove the two hypotheses from the cascade structure. The sign hypothesis
is the selection-weighted mean reversion of Prop.~34 \cite{NOTE} (LP
duality route sketched); the competition hypothesis ties the gap scale to
the increment scale --- the ultrametric spacing law (Lemma~26) and the
fine-end saturation (Prop.~23: $\CV_1 \to 0.5136 > 0$). Both are measured
across all available depths; neither is proved. We list the falsification
risks.
\end{enumerate}

\paragraph{Labels.} As in the companions, every assertion carries one of
four labels. \emph{Proved}: complete mathematical proof in the text or a
cited companion. \emph{Verified}: finite exact or floating-point
computation, reproducible from stated parameters. \emph{Measured}:
statistical estimate from computed data. \emph{Open}: neither proved nor
finitely checkable. Section~\ref{sec:status} collects the ledger.

\section{Setup and notation}
\label{sec:setup}

$\bet := \log_2 3$. For $j \geq 1$, $[3^j] := \{ m \bmod 3^j : m \equiv 2
\pmod 3 \}$. Fix $k \geq 2$ and $\lam \in (1,2]$. The nontrivial-tree
Krasikov--Lagarias system $L_k(\lam)$ asks for $c : [3^k] \to \R_{>0}$
with, for every $m$,
\begin{align}
m \equiv 2 \ (\mathrm{mod}\ 9): &\qquad
c(m) \leq \lam^{-2} c(4m) + \lam^{\bet-2}\, \bar c\!\left(\tfrac{4m-2}{3}\right),
\label{eq:L1}\\
m \equiv 5 \ (\mathrm{mod}\ 9): &\qquad
c(m) \leq \lam^{-2} c(4m),
\label{eq:L2}\\
m \equiv 8 \ (\mathrm{mod}\ 9): &\qquad
c(m) \leq \lam^{-2} c(4m) + \lam^{\bet-1}\, \bar c\!\left(\tfrac{2m-1}{3}\right),
\label{eq:L3}
\end{align}
where $4m$ is reduced mod $3^k$, branch targets mod $3^{k-1}$, and
$\bar c(t) := \min\{ c(t), c(t + 3^{k-1}), c(t + 2 \cdot 3^{k-1}) \}$ is
the minimum over the three refinements of $t$ --- the \emph{sibling
triple}. Feasibility implies $\pi_1(x) \geq x^{\gam}$ with
$\gam = \log_2 \lam$ \cite{KL03}. The certificate data below refer to the
critical (Perron) certificates \texttt{cert\_k13} etc.\ of
\cite{Temp26}; $P$ denotes the $3$-adic scale of the difference being
tracked, as in \cite{Drift26}.

\subsection{Triple coordinates}

Consider a sibling triple with member levels $x = (x_1, x_2, x_3)$ ---
the logarithmic certificate values whose minimum the system reads ---
and let one application of the edge-linearized operator move member $i$
by the increment $\Delta_i$. Decompose
\begin{equation}
\Delta_i = t + \dl_i, \qquad
t := \tfrac13 \textstyle\sum_i \Delta_i, \qquad
\textstyle\sum_i \dl_i = 0,
\label{eq:coords}
\end{equation}
so $t$ is the \emph{common (coarse) move} of the triple and $\dl_i$ the
centered \emph{fine moves}. Let
\begin{equation}
g_i := x_i - \min_j x_j \;\geq\; 0
\label{eq:gaps}
\end{equation}
be the static within-triple \emph{gaps} (at least one $g_i = 0$), and
write $g_{(1)} \leq g_{(2)} \leq g_{(3)}$ for their order statistics
($g_{(1)} = 0$; $g_{(2)}$ is the \emph{second gap}, the margin by which
the runner-up trails the argmin).

Statistics are taken over the population of sibling triples of a
certificate at a given scale $P$: $\Var(t)$ and $\cov(t, \cdot)$ are over
triples; $\Var(\dl) := \E[\dl_J^2]$ is the pooled fine variance, $J$ a
uniformly random member label independent of the triple (note
$\E[\dl_J] = 0$ and $\cov(t, \dl_J) = \E[t \,\E(\dl_J \mid
\mathrm{triple})] = 0$, both exact); $\sigma_\dl := \Var(\dl)^{1/2}$.

\begin{definition}[Attenuation constant; Def.~10 of \cite{Drift26}]
\label{def:kappa}
$\displaystyle \kap(P;k) := \frac{\mathrm{std}\big(\min_i (x_i + \Delta_i)
- \min_i x_i\big)}{\mathrm{std}(\Delta_J)}$,
the ratio of the standard deviation of the triple-minimum increment to
that of the pooled member increments.
\end{definition}

\section{The $S$-formula}
\label{sec:S}

\begin{lemma}[$S$-formula]
\label{lem:S}
\textup{(Proved.)}
For any triple, with the coordinates \eqref{eq:coords}--\eqref{eq:gaps},
\begin{equation}
\min_i (x_i + \Delta_i) \;-\; \min_i x_i \;=\; t + S,
\qquad
S := \min_i \,(\dl_i + g_i).
\label{eq:S}
\end{equation}
\end{lemma}

\begin{proof}
$\min_i (x_i + t + \dl_i)
= t + \min_i \big( \min_j x_j + g_i + \dl_i \big)
= \min_j x_j + t + \min_i (g_i + \dl_i)$.
\end{proof}

\emph{Verified:} identity \eqref{eq:S} holds to $1.8 \times 10^{-15}$
(maximum absolute residual, i.e.\ floating-point roundoff) on $150{,}000$
certificate triples.

\begin{corollary}[Slope exactly one; all attenuation in $S$]
\label{cor:slope}
\textup{(Proved.)}
The triple minimum passes the common move with slope exactly $1$: if
$\dl = 0$ then $S = \min_i g_i = 0$ and the min-increment equals $t$.
This recovers the directional low-pass, Prop.~34b of \cite{NOTE}
($\min(x + c\mathbf{1}) = \min(x) + c$), as the $\dl = 0$ case of
Lemma~\ref{lem:S}. In general the min-increment deviates from the common
move by exactly the scalar $S$, which depends only on the fine moves and
the gaps; moreover
\[
\min_i \dl_i \;\leq\; S \;\leq\; \dl_{i_0},
\qquad i_0 := \operatorname*{argmin}_i x_i ,
\]
the upper bound being the old-argmin half of the sandwich inequality of
Prop.~33 \cite{NOTE}.
\end{corollary}

\begin{proof}
$S \leq \dl_{i_0} + g_{i_0} = \dl_{i_0}$ since $g_{i_0} = 0$;
$S \geq \min_i \dl_i$ since $g_i \geq 0$.
\end{proof}

\begin{remark}[Relation to the clipping decomposition]
Prop.~33 of \cite{NOTE} decomposes the min-increment
as $\bar D + R$, member mean plus min-correction, and finds the
attenuation concentrated in the covariance $\cov(\bar D, R) < 0$. The coordinates $(t, S)$ of
Lemma~\ref{lem:S} are the same budget in a sharper basis: $t = \bar D$
and $S = R$, but now $S$ has a \emph{closed form} --- a minimum of fine
moves shifted by gaps --- rather than being defined as a residual. Every
structural question about the attenuation becomes a question about the
scalar random variable $S = \min_i(\dl_i + g_i)$.
\end{remark}

\section{The exact decomposition of $\kap$}
\label{sec:decomp}

\begin{corollary}[Exact $\kap$ decomposition]
\label{cor:decomp}
\textup{(Proved; verified exactly on certificate data.)}
\begin{equation}
\kap^2 \;=\;
\frac{\Var(t) + 2\cov(t,S) + \Var(S)}{\Var(t) + \Var(\dl)} ,
\label{eq:decomp}
\end{equation}
and consequently
\begin{equation}
\kap < 1
\quad\Longleftrightarrow\quad
2\cov(t,S) + \Var(S) \;<\; \Var(\dl) .
\label{eq:criterion}
\end{equation}
\end{corollary}

\begin{proof}
Numerator: by Lemma~\ref{lem:S} the min-increment is $t + S$, and
$\Var(t+S) = \Var(t) + 2\cov(t,S) + \Var(S)$. Denominator: with $J$ a
uniform member label independent of the triple,
$\Var(\Delta_J) = \Var(t + \dl_J) = \Var(t) + \Var(\dl)$, the cross term
vanishing exactly because $\E[\dl_J \mid \mathrm{triple}] = 0$.
Divide. The equivalence \eqref{eq:criterion} is numerator minus
denominator.
\end{proof}

\emph{Verified:} on \texttt{cert\_k13}, the right side of
\eqref{eq:decomp} equals the directly computed $\kap^2$ to all displayed
digits: $0.7901 = 0.7901$ at $P = 4$ and $0.7431 = 0.7431$ at $P = 6$
(i.e.\ $\kap = 0.889$, $0.862$, matching the independently measured
attenuation series $0.908 \to 0.841$ of \cite{Drift26} at those scales).

The decomposition splits the Open Lemma into exactly two mechanisms:
\begin{itemize}
\item[(M1)] \emph{Mean reversion}: $\cov(t,S) \leq 0$ --- the selected
fine move anti-correlates with the common move;
\item[(M2)] \emph{Min-variance contraction}: $\Var(S) < \Var(\dl)$ ---
the minimum over competing candidates is less variable than a single
member's fine move.
\end{itemize}
Either mechanism alone, if quantitatively uniform, would suffice; the
data (next section) shows \emph{both} present at every measured scale
and depth.

\section{Measured inputs}
\label{sec:measured}

All quantities on \texttt{cert\_k13}; covariances in normalized units
$V := \Var(t) + \Var(\dl) = 1$.

\begin{center}
\begin{tabular}{l c c}
\hline
quantity & $P = 4$ & $P = 6$ \\
\hline
$\kap^2$ (both sides of \eqref{eq:decomp}) & $0.7901$ & $0.7431$ \\
$\Var(S)/\Var(\dl)$ & $0.639$ & $0.499$ \\
$\cov(t,S)$ & $-0.0799$ & $-0.0517$ \\
$\Pr(S \text{ selects a non-argmin member})$ & $0.658$ & $0.665$ \\
median $g_{(2)} / \sigma_\dl$ & $0.130$ & $0.127$ \\
implied $\Var(\dl)/V$ (from \eqref{eq:decomp}) & $0.139$ & $0.306$ \\
\hline
\end{tabular}
\end{center}

\noindent Reading of the table (all four rows are \emph{measured}; the
last row is arithmetic from the rows above and closes the budget):

\begin{itemize}
\item \emph{Both mechanisms are active with the right sign.}
$\cov(t,S) < 0$ at both scales (M1), and $\Var(S)/\Var(\dl) < 1$ at both
scales (M2). The attenuation budget
$1 - \kap^2 = -2\cov(t,S) + \big(1 - \tfrac{\Var(S)}{\Var(\dl)}\big)
\tfrac{\Var(\dl)}{V}$
splits as $0.1598 + 0.0501$ at $P = 4$ (covariance-dominated) and
$0.1034 + 0.1535$ at $P = 6$ (contraction-dominated): the two mechanisms
trade weight with scale while their sum grows.
\item \emph{The competition is intense.} The median second gap is only
$0.13\,\sigma_\dl$: the runner-up typically trails the argmin by an
eighth of a fine-move standard deviation, so a single application
routinely reorders the triple. Consistently, $S$ selects a member other
than the static argmin with probability $0.658$--$0.665$ --- within one
percent of the value $2/3$ that full member-exchangeability with
gap-blind selection would predict. This is the data's vote for the
exchangeability hypothesis of Lemma~\ref{lem:contraction}(d) below.
\item \emph{The fine fraction grows with scale.} $\Var(\dl)/V$ rises
from $0.14$ to $0.31$ between $P = 4$ and $P = 6$: deep scales are
fine-dominated, which is where mechanism (M2) has the most leverage ---
consistent with the deepening of the measured attenuation
($\kap = 0.908 \to 0.841$ across $P$, \cite{Drift26}).
\end{itemize}

\section{The min-variance contraction}
\label{sec:contraction}

Mechanism (M2) asks when a minimum of competing random candidates is
\emph{less} variable than a single candidate. This is not automatic ---
a minimum of comonotone candidates just tracks the lowest one, with no
contraction --- so the content must lie in \emph{competition}. The
following lemma isolates exactly that: an identity with no hypotheses
whose deficit term is nonzero precisely when the \emph{order swaps},
plus a quantitative bound under hypotheses we state explicitly.

For a real random variable $W$ write $W^+ = \max(W, 0)$,
$W^- = \max(-W, 0)$.

\begin{lemma}[Min-variance contraction]
\label{lem:contraction}
\textup{(Proved; hypotheses as stated per part.)}
Let $A, B$ be square-integrable random variables on a common probability
space, $W := A - B$, $m := \min(A,B)$, $M := \max(A,B)$, and
$\sigma^2 := \max(\Var A, \Var B)$.
\begin{itemize}
\item[(a)] \emph{(Exact identity; no hypotheses.)}
\begin{equation}
\Var(m) + \Var(M) \;=\; \Var(A) + \Var(B) \;-\; 2\,\E[W^+]\,\E[W^-].
\label{eq:pairid}
\end{equation}
The deficit $2\E[W^+]\E[W^-]$ is strictly positive if and only if
$\Pr(A > B) > 0$ and $\Pr(B > A) > 0$: variance is destroyed exactly
when the order swaps.
\item[(b)] \emph{(Splitting min from max.)} With
$\Delta_{\mathrm{skew}} := \Var(m) - \Var(M)$,
\[
\Var(m) \;=\; \tfrac12\big(\Var A + \Var B\big) \;-\; \E[W^+]\,\E[W^-]
\;+\; \tfrac12 \Delta_{\mathrm{skew}} .
\]
If the law of $(A,B)$ is invariant under the anti-diagonal reflection
$(A,B) \mapsto (s - B,\, s - A)$ for some constant $s$, then
$\Delta_{\mathrm{skew}} = 0$ and
\[
\Var(\min(A,B)) \;=\; \tfrac12\big(\Var A + \Var B\big)
\;-\; \E[W^+]\,\E[W^-] \;\leq\; \sigma^2 - \E[W^+]\,\E[W^-].
\]
\item[(c)] \emph{(Quantitative competition bound.)} If, in addition to
the reflection hypothesis of (b), there are $p_0 > 0$, $c > 0$ with
\[
\Pr(W \geq c\,\sigma) \;\geq\; p_0
\qquad\text{and}\qquad
\Pr(-W \geq c\,\sigma) \;\geq\; p_0
\]
(two-sided order-swap frequency), then
\[
\Var(\min(A,B)) \;\leq\; \big(1 - \varepsilon(p_0, c)\big)\,\sigma^2,
\qquad
\varepsilon(p_0, c) := p_0^2\, c^2 \;>\; 0 .
\]
\item[(d)] \emph{(Exchangeable triple.)} Let $(Y_1, Y_2, Y_3)$ be
exchangeable with common variance $\sigma^2$ and mean $\mu$, order
statistics $Y_{(1)} \leq Y_{(2)} \leq Y_{(3)}$, and set
$a_i := \E[Y_{(i)}] - \mu$. Then, with no further hypotheses,
\[
\textstyle\sum_i \Var(Y_{(i)}) \;=\; 3\sigma^2 - \sum_i a_i^2 ,
\]
and if additionally the vector is centrally symmetric
($2\mu\mathbf{1} - Y \overset{d}{=} Y$), then $\Var(Y_{(1)}) =
\Var(Y_{(3)})$, $a_2 = 0$, and
\[
\Var(\min_i Y_i) \;=\; \tfrac12\Big( 3\sigma^2
\;-\; \tfrac12\big(\E[\,Y_{(3)} - Y_{(1)}\,]\big)^2
\;-\; \Var(Y_{(2)}) \Big) .
\]
In particular $\Var(\min) < \sigma^2$ whenever
$\big(\E[\mathrm{range}]\big)^2 + 2\Var(\mathrm{median}) > 2\sigma^2$.
\end{itemize}
\end{lemma}

\begin{proof}
(a) Pointwise $\{m, M\} = \{A, B\}$, so $\E[m^2 + M^2] = \E[A^2 + B^2]$.
From $m = \tfrac12(A+B) - \tfrac12|W|$ and $M = \tfrac12(A+B) +
\tfrac12|W|$, with $\mu := \tfrac12(\E A + \E B)$, $d := \E|W|$,
$b := \E W$:
\[
(\E m)^2 + (\E M)^2 = 2\mu^2 + \tfrac{d^2}{2},
\qquad
(\E A)^2 + (\E B)^2 = 2\mu^2 + \tfrac{b^2}{2}.
\]
Subtracting, $\Var m + \Var M = \Var A + \Var B - \tfrac12(d^2 - b^2)$.
Since $|W| = W^+ + W^-$ and $W = W^+ - W^-$,
$d - b = 2\E[W^-]$ and $d + b = 2\E[W^+]$, so
$d^2 - b^2 = 4\,\E[W^+]\E[W^-]$, giving \eqref{eq:pairid}. The deficit is
positive iff both factors are, i.e.\ iff the order swaps both ways with
positive probability.

(b) The first display is \eqref{eq:pairid} rearranged. Under the
reflection, $\min(s - B, s - A) = s - \max(A, B)$, so
$m \overset{d}{=} s - M$ and $\Var m = \Var M$.

(c) $\E[W^+] \geq c\sigma \Pr(W \geq c\sigma) \geq p_0 c \sigma$ and
symmetrically for $W^-$; insert into (b) and use
$\tfrac12(\Var A + \Var B) \leq \sigma^2$.

(d) Pointwise $\sum_i Y_{(i)}^2 = \sum_i Y_i^2$, so the second moments
agree; exchangeability gives $\E Y_i = \mu$ for all $i$, and
$\sum_i a_i = 0$, so $\sum_i (\E Y_{(i)})^2 = 3\mu^2 + \sum_i a_i^2$,
proving the first display. Central symmetry maps
$Y_{(1)} \mapsto 2\mu - Y_{(3)}$ and $Y_{(2)} \mapsto 2\mu - Y_{(2)}$,
so $\Var Y_{(1)} = \Var Y_{(3)}$, $a_2 = 0$, and
$a_3 = -a_1 = \tfrac12 \E[Y_{(3)} - Y_{(1)}]$, whence
$\sum a_i^2 = \tfrac12 (\E[\mathrm{range}])^2$; solve the first display
for $\Var(Y_{(1)})$.
\end{proof}

\begin{remark}[Near-ties plus dispersion imply order swaps]
\label{rem:nearties}
The intuitive competition event --- ``at least two candidates lie within
$c_0 \sigma_\dl$ of each other'' --- does not by itself produce
contraction: two comonotone candidates a hair apart never swap, and part
(a) then gives deficit zero. What near-ties buy is \emph{leverage} for
the fine-move dispersion. Concretely, for the two lowest-gap members $a$
(argmin, $g_a = 0$) and $b$ (runner-up, $g_b = g_{(2)}$),
$W = Y_a - Y_b = (\dl_a - \dl_b) - g_{(2)}$; on the joint event
$\{ g_{(2)} \leq c_0 \sigma_\dl \} \cap
\{ \dl_a - \dl_b \geq c_1 \sigma_\dl \}$ with $c_1 > c_0$ one has
$W \geq (c_1 - c_0)\sigma_\dl$, while
$\{-W \geq c_1 \sigma_\dl\} \supseteq \{\dl_b - \dl_a \geq c_1\sigma_\dl\}
\cap \{g_{(2)} \geq 0\}$ costs nothing. So: \emph{small static gaps}
(measured: median $g_{(2)} = 0.13\,\sigma_\dl$) \emph{plus two-sided
anti-concentration of the fine-move differences} (non-comonotonicity)
deliver the two-sided swap hypothesis of part (c) with $c = c_1 - c_0$.
The near-tie condition and the swap condition are the static and dynamic
faces of the same competition.
\end{remark}

\begin{remark}[Which hypotheses does the data support?]
\label{rem:hypotheses}
Honesty requires separating the three hypotheses used above.
\begin{itemize}
\item \emph{Exchangeability} (part (d), and implicitly the ``no
distinguished member'' reading of (c)): supported. Under member-\linebreak
exchangeability with gaps small relative to fine moves, the selected
member is nearly uniform over the triple, predicting
$\Pr(\text{non-argmin selection}) = 2/3 = 0.667$; measured $0.658$ and
$0.665$ (Section~\ref{sec:measured}).
\item \emph{Competition / order swaps}: supported. Median
$g_{(2)}/\sigma_\dl = 0.130$, $0.127$ --- the near-tie regime of
Remark~\ref{rem:nearties} --- and the $0.66$ non-argmin selection rate
is direct evidence that the order swaps at nearly the exchangeable
ceiling.
\item \emph{Reflection / central symmetry} (the hypothesis that kills
the skew term $\Delta_{\mathrm{skew}}$): \emph{untested}. We have not
separately measured $\Var(\min) - \Var(\max)$ on the certificate
triples; a persistent positive skew would weaken part (c)'s bound by
$\tfrac12\Delta_{\mathrm{skew}}$. What \emph{is} directly verified is
the conclusion itself: $\Var(S)/\Var(\dl) = 0.639$, $0.499 < 1$ --- so
the package (swap frequency $+$ approximate symmetry) is consistent ex
post, but the symmetry hypothesis is a genuine assumption of the clean
statement, not a theorem about the cascade. We flag it as such in
Theorem~\ref{thm:main} and Section~\ref{sec:gap}.
\end{itemize}
Note also a bookkeeping point: parts (a)--(d) control $\Var(\min Y)$
relative to $\Var(Y_i) = \Var(\dl_i + g_i)$, which exceeds $\Var(\dl)$
when the gaps fluctuate across triples. Converting the contraction into
the form needed by \eqref{eq:criterion}, $\Var(S) \leq (1 -
\varepsilon_0)\Var(\dl)$, therefore consumes part of the deficit to pay
for the gap variance --- which is affordable precisely in the near-tie
regime where gaps are small on the $\sigma_\dl$ scale. This is where the
gap-to-increment scale ratio enters irreducibly; see
Section~\ref{sec:gap}.
\end{remark}

\section{The conditional stability theorem}
\label{sec:main}

\begin{theorem}[Conditional uniform attenuation]
\label{thm:main}
\textup{(Proved as stated; the hypotheses are the open content.)}
Suppose there exist constants $\varepsilon_0 > 0$ and $\varphi_0 > 0$
such that, for every scale $P$ and every depth $k$, in the stationary
(Perron) state of the critical Krasikov--Lagarias system:
\begin{itemize}
\item[(i)] \emph{(Mean-reversion sign)} $\cov(t, S) \leq 0$;
\item[(ii)] \emph{(Uniform competition)}
$\Var(S) \leq (1 - \varepsilon_0)\,\Var(\dl)$ \ and \
$\Var(\dl) \geq \varphi_0 \big( \Var(t) + \Var(\dl) \big)$.
\end{itemize}
Then
\[
\kap(P; k)^2 \;\leq\; 1 - \varepsilon_0\,\varphi_0 \;=:\; \kap_{\max}^2
\;<\; 1
\qquad \text{for all } P, k :
\]
the Open Lemma of \cite{Drift26, Gamma26} holds with
$\kap_{\max} = (1 - \varepsilon_0 \varphi_0)^{1/2}$, and consequently, by
the conditional theorem of \cite{Gamma26} (whose remaining five links are
proved there), $\gam_k \to 1$ and
\[
\pi_1(x) \;\geq\; x^{1 - \varepsilon}
\qquad \text{for every } \varepsilon > 0 .
\]
\end{theorem}

\begin{proof}
By Corollary~\ref{cor:decomp} and (i), then (ii),
\[
\kap^2 \leq \frac{\Var(t) + \Var(S)}{\Var(t) + \Var(\dl)}
\leq \frac{\Var(t) + (1 - \varepsilon_0)\Var(\dl)}{\Var(t) + \Var(\dl)}
= 1 - \varepsilon_0 \,\frac{\Var(\dl)}{\Var(t) + \Var(\dl)}
\leq 1 - \varepsilon_0 \varphi_0 .
\]
\end{proof}

\begin{remark}[Measured values of the constants]
The hypotheses hold with room at every measured point:
$\cov(t,S) = -0.0799$, $-0.0517 < 0$;
$\varepsilon_0 = 1 - 0.639 = 0.361$ at the weakest measured scale;
$\varphi_0 = 0.139$ at the coarsest measured scale (and growing with
$P$). These give $\kap_{\max} = (1 - 0.361 \times 0.139)^{1/2} = 0.975$
--- a valid but crude bound, since the proof spends hypothesis (i) only
on its sign. Keeping the measured covariance yields the measured
$\kap = 0.889$ at $P = 4$; the theorem's job is uniformity, not
sharpness. Note also that Lemma~\ref{lem:contraction}(c), under its
stated hypotheses with $(p_0, c)$ uniform in depth and the gap-variance
bookkeeping of Remark~\ref{rem:hypotheses}, is precisely a mechanism
that would deliver hypothesis (ii); and Prop.~34's binding rigidity
\cite{NOTE} is the mechanism proposed for (i). Neither delivery is yet a
proof.
\end{remark}

\section{The honest gap}
\label{sec:gap}

Theorem~\ref{thm:main} is an exact reduction, not a resolution: the Open
Lemma is now equivalent (via the verified identity \eqref{eq:decomp}) to
proving hypotheses (i) and (ii) \emph{from the cascade structure}. We
record what each requires, what is measured, and what would falsify the
program.

\subsection{Proving (i): selection-weighted mean reversion}

$\cov(t, S) \leq 0$ says: when the triple's common move is up, the
selected member's fine move is preferentially down, and vice versa. This
is exactly the phenomenon isolated in Prop.~34 of \cite{NOTE}: the
antisymmetric min-correction ($\pm 0.026$) is carried equally by switch
and no-switch events, so the mechanism is not clipping-at-switches but
\emph{binding-constraint rigidity} --- the argmin member is the binding
entry of the rows that read it, and binding entries, pinned by the
equation network through complementary slackness, co-move less than the
triple average. The proof route (sketched in \cite{Drift26},
Sect.~``The LP-duality route'') is:
\begin{enumerate}
\item at the critical $\lam_k$ the certificate is an LP optimum, and
each triple's argmin carries a strictly positive dual weight;
\item show the dual weights concentrate on binding entries with pinning
strength bounded below uniformly, using the mass-conservation identity
(normalizing total dual mass) and the balance identity (preventing dual
mass from draining across the scale hierarchy), both proved in
\cite{Drift26};
\item sensitivity analysis then bounds the binding member's response
slope below the slack members', giving $\cov(t, S) \leq 0$ in the
stationary state.
\end{enumerate}
Step 2 is the missing arithmetic input; nothing here is circular with
the present paper, but nothing here is proved either.

\subsection{Proving (ii): the gap scale versus the increment scale}

Uniform competition requires that the static gaps do not outgrow the
fine-move dispersion: the near-tie ratio $g_{(2)}/\sigma_\dl$ must stay
bounded (measured: median $0.130$, $0.127$ at $P = 4, 6$) and the fine
fraction $\Var(\dl)/V$ must stay bounded below (measured: $0.139$,
$0.306$ and growing with scale). Two measured laws of \cite{NOTE} are
the natural sources:
\begin{itemize}
\item \emph{The ultrametric spacing law} (Lemma~26): the triple-$\CV$ of
the certificate at class-spacing $s$ depends to leading order only on
$v_3(s)$, with an approximate two-term closure. This is exactly a
statement tying the gap distribution at each scale to the certificate's
own $\CV$ profile --- i.e.\ tying the gap scale to the increment scale.
Measured, with an emergent-coefficient caveat; not proved.
\item \emph{Fine-end saturation} (Prop.~23): $\CV_1(k) = 0.5136 -
0.337\,(0.910)^k$ to residual $10^{-6}$ across nine depths --- the
intra-triple variation converges to a \emph{nonzero} limit. This is
what keeps $\sigma_\dl$ (and with it the fine fraction $\varphi_0$)
alive as $k \to \infty$; its failure ($\CV_1 \to 0$) would kill
competition and push $\kap \to 1$. Measured with striking precision;
not proved.
\end{itemize}

\subsection{What is measured across depth}

The program's empirical safety margin is Prop.~35 of \cite{NOTE}:
$\kap < 0.95$ at every scale and depth, with the deep-scale value flat,
$\kap_{\mathrm{deep}} = 0.841, 0.840, 0.837, 0.838$ at
$k = 13, 15, 17, 19$, i.e.\ $\kap_\infty = 0.839 \pm 0.002$. And by the
scalar closure (Prop.~47 of \cite{NOTE}), the variance budget of the
exact transport identity closes to $0.06\%$ with $\kap(P)$ as the only
nonlinear unknown: the entire cascade is governed by a one-dimensional
map whose measured trajectory descends ($0.917 \to 0.899$ across $P$,
deep limit $0.839$). The stability program of the title is precisely:
prove that this one-dimensional fixed point sits below $1$ --- which
Theorem~\ref{thm:main} reduces to the sign (i) and the competition (ii).

\subsection{Falsification risks}

Stated symmetrically, as in the companions:
\begin{enumerate}
\item \emph{Sign creep.} $\cov(t,S)$ could drift to $0$ or turn positive
at depths beyond measurement. The measured normalized value shrinks from
$-0.0799$ ($P{=}4$) to $-0.0517$ ($P{=}6$); on the present data the
contraction term more than compensates (total attenuation grows), but a
proof of (i) must rule out asymptotic sign loss, and the trend is a
genuine caution.
\item \emph{Competition dilution.} The gap scale could grow relative to
$\sigma_\dl$ at deep scales or depths ($g_{(2)}/\sigma_\dl$ measured
flat at $0.13$, but only to $P = 6$, $k = 19$), or the fine fraction
could collapse ($\varphi_0 \to 0$) if Prop.~23's saturation is
preasymptotic.
\item \emph{Symmetry breakdown.} The clean contraction bound assumes
reflection symmetry; a persistent skew $\Delta_{\mathrm{skew}} > 0$
would eat into the deficit $\E[W^+]\E[W^-]$. This hypothesis is
currently untested on certificate data (Remark~\ref{rem:hypotheses}) and
is the cheapest of the risks to retire: measure
$\Var(\min) - \Var(\max)$ across $P$ and $k$.
\item \emph{Exchangeability breakdown.} The non-argmin selection rate
could drift from $2/3$ at large $k$, signaling a distinguished member
and invalidating the symmetric reading of the competition bound.
\item \emph{Preasymptotic flatness.} The depth-flatness of
$\kap_{\mathrm{deep}}$ over $k = 13$--$19$ spans a factor $3^6$ in
modulus but is still a finite window; all conclusions drawn from it are
measured, not proved, and are labeled so.
\end{enumerate}

\section{Status ledger}
\label{sec:status}

\begin{center}
\begin{tabular}{p{0.22\textwidth} p{0.71\textwidth}}
\hline
\textbf{Proved} &
$S$-formula \eqref{eq:S} and slope-one corollary
(Lemma~\ref{lem:S}, Cor.~\ref{cor:slope}), recovering Prop.~34b;
exact decomposition \eqref{eq:decomp} and criterion \eqref{eq:criterion}
(Cor.~\ref{cor:decomp});
min-variance identities and contraction bounds under stated hypotheses
(Lemma~\ref{lem:contraction}(a)--(d));
conditional uniform attenuation (Thm.~\ref{thm:main});
the five other links of the $\gam \to 1$ chain \cite{Temp26, Drift26,
Gamma26}. \\
\hline
\textbf{Verified} &
$S$-formula to $1.8 \times 10^{-15}$ on $150{,}000$ certificate triples;
decomposition closure $0.7901 = 0.7901$ ($P{=}4$), $0.7431$ ($P{=}6$),
$k = 13$. \\
\hline
\textbf{Measured} &
$\Var(S)/\Var(\dl) = 0.639$, $0.499$;
$\cov(t,S) = -0.0799$, $-0.0517$;
non-argmin selection $0.658$, $0.665$ (vs.\ $2/3$ exchangeable);
median $g_{(2)}/\sigma_\dl = 0.130$, $0.127$;
implied fine fraction $0.139$, $0.306$;
$\kap$ flat at $0.839 \pm 0.002$ over $k = 13$--$19$ (Prop.~35);
$\CV_1 \to 0.5136 > 0$ (Prop.~23); spacing law (Lemma~26);
scalar closure to $0.06\%$ (Prop.~47) \cite{NOTE}. \\
\hline
\textbf{Open} &
Hypothesis (i): $\cov(t,S) \leq 0$ uniformly, from binding-constraint
rigidity (Prop.~34; LP-duality route, step 2 missing).
Hypothesis (ii): uniform competition --- gap scale tied to increment
scale (Lemma~26) and fine-end saturation (Prop.~23), neither proved.
Reflection-symmetry hypothesis of
Lemma~\ref{lem:contraction}(b)--(c): untested on certificate data. \\
\hline
\end{tabular}
\end{center}

\subsection*{Acknowledgment of method}
The computations, fits, and drafts were produced in an extended
human--AI collaboration; every proved statement above has a
self-contained proof in the text or in a cited companion, independent of
that provenance.

\begin{thebibliography}{9}

\bibitem{KL03}
I.~Krasikov and J.~C.~Lagarias,
\emph{Bounds for the 3x+1 problem using difference inequalities},
Acta Arith. \textbf{109} (2003), no.~3, 237--258. arXiv:math/0205002.

\bibitem{Lag85}
J.~C.~Lagarias,
\emph{The 3x+1 problem and its generalizations},
Amer. Math. Monthly \textbf{92} (1985), 3--23.

\bibitem{Tao19}
T.~Tao,
\emph{Almost all orbits of the Collatz map attain almost bounded values},
Forum Math. Pi \textbf{10} (2022), e12. arXiv:1909.03562 (2019).

\bibitem{Temp26}
M.~de Jong (with Jengo, AI),
\emph{Krasikov--Lagarias certificates for the $3x+1$ problem are tempered
roulette measures}, companion paper, July 2026.

\bibitem{Drift26}
M.~de Jong (with Jengo, AI),
\emph{The drift of the Collatz cascade: balance identities, the directional
low-pass, and the attenuation constant}, companion paper, July 2026.

\bibitem{Gamma26}
M.~de Jong (with Jengo, AI),
\emph{Towards density $x^{1-\varepsilon}$ for the $3x+1$ problem: a program
with one open lemma}, companion paper, July 2026.

\bibitem{NOTE}
M.~de Jong (with Jengo, AI),
\emph{Six results on the Collatz map in family/pair coordinates}, research
note, July 2026 (Propositions 23, 33, 34, 34b, 35, 47; Lemma 26).

\end{thebibliography}

\end{document}
